\subsection{KV Placement Planner: Blocks and Layers}
\label{sec:kv-placement}

Planner (2) describes a request's Peer protection as the hottest $K_l$
sealed blocks of each of the first $L$ layers. Calibration determines
the hot-set sizes; online rankings determine their identities; a
transfer--compute model determines the front-layer prefix.

\paragraph{Protection unit and replay endpoint.}
The unit is a \emph{(layer, block)} containing all KV heads. AegisKV uses
64-token blocks for the profiled selectors, aligning counting, sealing,
and copying with their block granularity. Larger units bundle cold and
hot tokens; smaller units increase object and metadata counts.
At the same 20\% byte budget, tying each token span to all layers covers
only 53--59\% of selection events in the 512-query traces, versus 66\%
for MiniCPM4-8B and over 90\% for NOSA-1B with layer-block selection.

A full block is sealed and becomes eligible for copying. During decode,
the unsealed tail is shorter than one block; the analysis conservatively
uses 64 tokens (about 32 in expectation under a uniform failure phase).
Sealed but uncommitted KV also belongs to the replay suffix, determined
by the commit frontier. Replay forwards the already generated token IDs
to rebuild missing KV; it neither resamples those tokens nor recomputes
an available historical prefix. Our recovery endpoint is completion of
this known-tail KV reconstruction.

\paragraph{Size the hot set by distinct needed blocks.}
For layer $l$, replaying $n$ tail tokens needs
\begin{equation}
 U_l(n)=\bigcup_{q=1}^{n}\ \bigcup_g S_{l,q,g},
\end{equation}
where $S_{l,q,g}$ is the selector output for query $q$ and selection
group $g$. A repeatedly selected block appears once in $U_l$.
At 32K history, the 64-query union averages 37\% of history bytes for
NOSA-1B and 68\% for MiniCPM4-8B; for 512 queries these rise to 58\%
and 90\%. Thus, event-frequency concentration cannot determine required
recovery bytes.

A table indexed by model, layer, context bin, and tail length estimates
$h_l=|U_l|/B_l$, where $B_l$ is the number of historical blocks in layer
$l$. It yields $K_l=\lceil h_l B_l\rceil$, validated to include the
selector's mandatory historical blocks within that budget.
In the current three-document calibration, a table from other documents
gives a modeled demand floor of 20.3\,ms, compared with 19.9\,ms from
window-specific estimation. Adjacent-window union fractions change by
0.03 on average (P90: 0.07). Layer specificity matters: fractions vary
from 0.44 to 0.87 across MiniCPM4 layers; replacing them with one constant
raises this floor to 23.5\,ms. Context matters too: the mean calibrated
fraction is 0.89 at 16K versus 0.66 at 32K. The Adapter can report union
statistics for slow, cross-request profile refreshes rather than
resizing every request's hot set online.

\paragraph{Track identities online.}
The Adapter reads indices already produced by native attention kernels
and accumulates selection counts in a layer-by-block table. It assumes
neither identical selectors nor a fixed selection budget across models.
In the profiled MiniCPM4 configuration, top-64 selection plus a
2,048-token local window selects at most 96 blocks per query per layer;
NOSA-1B selects 64 including mandatory blocks. With 32 layers, two KV
heads, and 32K history, MiniCPM4 needs at most 6,144 counter increments
per decoding step and a 32\,KiB table of 16-bit counters. Deduplicating
indices within each group and resetting every 512 steps bounds a
counter by 1,024. These are operation and storage counts, not measured
runtime overhead.

At each ranking window, the Planner prioritizes mandatory sealed blocks
and fills the remaining $K_l$ slots by decreasing count, breaking ties
by recency. An old optional block becomes replaceable only after two
successive windows outside the top set. Replacements use existing slots,
so hysteresis does not inflate the $K_l$ budget. Newly selected objects
trigger incremental copies; removed objects retain DRAM coverage.
Figure~\ref{fig:design-coverage} extends the budget comparison in
Figure~\ref{fig:insight-retry-kv}. At 40\% historical bytes, past-frequency
ranking covers 79.2\%/77.8\%/53.6\% of distinct needed KV for
NOSA-1B/NOSA-3B/MiniCPM4, versus 59.7\%/59.7\%/50.9\% for recency.
This isolates the ranking signal with a fixed per-layer budget; it does
not evaluate the complete admission policy or two-window hysteresis.

\begin{figure}[t]
  \centering
  \includegraphics[width=\columnwidth]{figures/design-insights/retry-kv-budget.pdf}
  \caption{Budget sensitivity of the same held-out KV forecast as
    Figure~\ref{fig:insight-retry-kv}. Both policies include mandatory sink
    and local-window blocks within the plotted budget. Coverage sums
    distinct covered/needed historical layer-blocks per document, then
    averages over three documents per model. Solid/filled curves use
    past frequency; dashed/hollow curves use recency. Each model has one
    native capture per document; budget points reuse those captures.}
  \label{fig:design-coverage}
  \Description{Past-window hotness covers needed KV blocks earlier than
    recency, with a larger difference for NOSA than for MiniCPM4.}
\end{figure}

\paragraph{Preload and demand fetch.}
The recovery loader preloads $K_l$ hot blocks for each layer. Selector
state, such as compressed keys or block summaries, must be available
before selecting that layer's exact demand. Once demand is known, misses
receive higher priority than speculative preloads; remaining cold KV is
filled in the background. Let $\widetilde u_l$ denote the fraction of
historical KV in the union of preload and demand-fetch sets. Replaying
64 queries at 32K gives mean $\widetilde u=0.44$ for NOSA-1B and 0.77
for MiniCPM4, about 5--10 percentage points above the ideal needed union.
At 16K these values are 0.69 and 0.96. They describe trace-derived
sets, not a measured DMA schedule. Dense attention approaches
$h=\widetilde u=1$ and reduces to full-layer loading.

\paragraph{Protect the front-layer preload backlog.}
Consider $N$ homogeneous layers, with mean preload fraction $h$,
effective fraction $\widetilde u$, per-layer tail compute time $C$,
and full-history transfer time $D=L_rs/\beta$.
Here $s$ is KV bytes per token per layer and $\beta$ is effective
bandwidth under representative concurrency. Preloads can be queued in
advance, whereas a layer's misses are released only when its selector
can determine demand. With an ideal preemptible copy stream,
\begin{align}
 W_{\mathrm{floor}}&=N(\widetilde u-h)D,\label{eq:demand-floor}\\
 W_0&=hD+(N-1)[hD-C]^++W_{\mathrm{floor}}.\label{eq:exposed-wait}
\end{align}
The first two terms are startup preload and the backlog that computation
cannot hide. The floor is demand-fetch wait under this \emph{fixed hot
set and release policy}. Increasing the number of Peer-protected layers
removes preload wait, not misses outside the hot set. Enlarging or
changing that set can lower the floor; it is not a bound on all placements.

Figure~\ref{fig:design-pipeline} illustrates the progression at 32K history
with the measured MiniCPM4 tail cost and $\beta=6.25$\,GB/s.
Waiting for all historical KV before computing exposes 172\,ms;
full-layer pipelining reduces this to 83\,ms, and hot preloads with
demand fetching to 44\,ms. Keeping the first seven layers' hot sets
(147\,MiB) in Peer HBM leaves 20\,ms of demand wait. These are modeled
timelines, not end-to-end fault-recovery measurements.

\begin{figure*}[t]
  \centering
  \includegraphics[width=\textwidth]{figures/design-paper/design-layer-pipeline.pdf}
  \caption{Calibrated single-request recovery illustration: MiniCPM4-8B,
    32 layers, 32K history (32\,MiB/layer), and a known 64-token tail.
    (a) Composes measured 4.70\,s history prefill with calibrated tail
    compute; it is not a measured one-shot 32K+64 run.
    (b--e) Use an expanded millisecond axis, $D=5.37$\,ms at
    6.25\,GB/s, and $C=2.87$\,ms from A100 measurements. Rows compare
    waiting for all history, full-layer pipelining, preload plus demand
    fetch, and the same loader with front-layer hot sets in Peer HBM.
    Waits are 172, 83, 44, and 20\,ms; totals include 92\,ms of tail
    compute and end at reconstructed tail KV, not cold-copy completion.
    The model uses homogeneous layer/document means and a preemptible
    copy stream; it excludes detection, queuing, and resource interference.}
  \label{fig:design-pipeline}
  \Description{A seconds-scale recomputation reference and four
    millisecond-scale timelines separate transfer, tail compute, and
    exposed KV wait. Peer-resident front-layer hot sets reduce preload
    wait while demand-fetch misses remain.}
\end{figure*}

The Peer byte target for eliminating the preload component is
\begin{equation}
 F^\star(r)=\beta\bigl(W_0-W_{\mathrm{floor}}\bigr)
 =\beta\bigl(hD+(N-1)[hD-C]^+\bigr).
 \label{eq:peer-byte-target}
\end{equation}
Let $b$ be bytes per layer-block. Choose the smallest front-layer prefix
$L$ whose hot-set bytes reach $F^\star$; the KV footprint is
$F_{\mathrm{KV}}(r)=\sum_{l=1}^{L}K_l b$.
Admission also counts selector-state bytes. If $hD\le C$, the first
layer's hot set suffices. Figure~\ref{fig:design-footprint} scans the same
causal pipeline: full-layer loading needs 16 layers/512\,MiB to eliminate
wait, the MiniCPM4 profile needs seven/147\,MiB to reach its demand floor,
and the NOSA-1B profile needs one/11\,MiB. The NOSA curve retains the
MiniCPM4 layer sizes and compute cost to isolate the access profile.
It is not a measured NOSA recovery latency.

\begin{figure}[t]
  \centering
  \includegraphics[width=\columnwidth]{figures/design-wait-vs-bytes.pdf}
  \caption{Modeled exposed wait versus front-layer hot-set residency.
    All curves use the 32-layer MiniCPM4 byte and compute configuration
    of Figure~\ref{fig:design-pipeline}; only $(h,\widetilde u)$ changes:
    $(1,1)$, $(0.66,0.77)$, and $(0.36,0.44)$.
    One step adds one layer's hot set. Triangles mark the first prefix
    reaching the corresponding fixed-hot-set demand floor.
    Footprints do not depend on slack.}
  \label{fig:design-footprint}
  \Description{Dense loading reaches zero wait at 512 MiB; MiniCPM4 and
    NOSA access profiles reach nonzero demand floors with smaller footprints.}
\end{figure}

The footprint depends on context, model profile, and bandwidth, not slack.
Within a fixed prefix length it varies approximately linearly with
context, with discrete changes when $L$ changes. The profiled 16K case
needs one layer/14\,MiB; extrapolating the 32K union fraction to 64K and
128K gives about 20 layers/843\,MiB and 25 layers/2.1\,GiB.
These larger-context footprints are extrapolations. At 25\,GB/s,
the 32K case needs one layer/21\,MiB. A 32-token rather than 64-token
tail reduces the illustrative footprint from 147 to 92\,MiB.
$L$ and the footprint are fixed within a placement window and recomputed
at its boundary. Equation~\ref{eq:peer-byte-target} is a planning
approximation for heterogeneous layers, not a guarantee that their
actual runtime reaches exactly the homogeneous floor.
